%--
\chapter{线性时间选择}

%----------------------------------------
\section{分组大小背后的微妙权衡}

在\alg{Select-WLinear}算法中，有一个“神奇”的数字“5”。
想要深入理解\alg{Select-WLinear} 算法，一个自然的疑问就是，为什么元素的分组大小应该是5？
或者说，其他的分组大小——例如分组大小为3、4、7、9等——行不行？

看起来分组的大小有无数种选项，但是我们可以经过简单的分析，将关注点聚焦到少数代表性选项，这可以帮助我们有效理解分组大小“5”背后的原理。
首先不难理解，偶数的分组大小可以被排除，因为每组有奇数个元素时，不仅中位数的定义更加简洁，而且更有利于保证最终划分的平衡性。
其次，分组大小的选择，以5为基准，有两个方向，一个方向是将分组变得更“细”，此时只有分组大小为3这一个选择。另一个方向是将分组变得更“粗”，这有无穷多种选择，但是我们可以对分组大小7做一个深入剖析，进而理解分组大小选择5的原理。所以我们下面对代表性的分组大小3和7，进行深入分析。

%----------
\subsubsection{分组大小为3的情况}

如果我们将\alg{Selection-WLinear}算法中分组的大小定为3，而算法的其他设计保持不变，此时算法的运作更接近于\alg{Selection-ELinear}中所有元素一字排开的情况。
考虑到\\ \alg{Selection-ELinear}算法的最坏情况时间复杂度$O(n^2)$，我们有理由担心算法的最坏情况时间复杂度是否能保持在$O(n)$的水平。
基于上面的考虑，我们试图去证明算法时间复杂度的一个下界，若能证明算法的时间复杂度严格超过了$O(n)$，则可以表明分组大小3不可行。

根据算法的划分方式，我们可以确定比$m^*$小的元素，至多有：
%--
\begin{align*}
    2 \times \ceil{ \frac{1}{2} \ceil{\frac{n}{3}} } \leq \frac{n}{3}+3
\end{align*}

\noindent 上述不等式左边的原理很简单，大致估算下来，所有分组中，有一半的小组，每个小组有2个元素，确定比$m^*$小。由于此时我们只需要估算一个上界，所以我们在除以2或者3的时候，都取上取整，在计算小组数量时，也未扣除$m^*$所在的小组。
%
与\alg{Selection-WLinear}算法的分析类似，此时最坏情况下，我们需要对至少$n-(\frac{n}{3}+3) = \frac{2}{3}n-3$个元素进行递归。
所以我们有算法的最坏情况时间复杂度满足：
%--
\begin{align*}
    W(n) \geq W\bkt{\ceil{\frac{n}{3}}} + W\bkt{\frac{2}{3}n-3} + \Omega(n)
\end{align*}

\noindent $W(n)$递归不等式右边的第一项，是对所有小组的中位数，选择$m^*$的代价；第二项是最坏情况下，对最多的元素进行递归选择的代价；第三项是元素划分的代价。
通过“guess and prove”的方式（参见\ref{SubSec:Guess-and-Prove}小节），我们不难证明$W(n) = \Omega(n\log n)$，具体证明留作习题。

上述分析表明，分组大小若是设为3，则$W(n)$的渐进增长率严格高于$O(n)$，所以分组大小为3的选项被\alg{Selection-WLinear}算法合理抛弃。

%----------
\subsubsection{分组大小为7的情况}

以分组大小5为基准，更大的取值代表着更精细的分组和更好地平衡性控制（当$n$充分大的时候），所以我们有理由相信算法的最坏情况代价仍然能保持$O(n)$的水平，而不会出现分组大小为3时，$W(n) = \Omega(n\log n)$的情况。基于这一合理猜测，我们来严格证明$W(n)=O(n)$。

考虑7个元素一小组的划分方式，则确定比$m^*$小的元素至少有：
%--
\begin{align*}
    4 \times \bkt{\ceil{\frac{1}{2}\ceil{\frac{n}{7}}}-2} \geq \frac{2}{7}n - 8
\end{align*}

\noindent 上述不等式左边的原理是，我们大致有一半约$\frac{n}{7}\times \frac{1}{2}$个小组，每个小组贡献4个元素（每组上方有3个元素，再加上组内的中位数）确定比$m^*$小。在处理取整的时候，我们为了确定得到一个下界，所以保守地去掉了$m^*$所在的小组，也去掉了最后一个可能并不完整的小组（由于$n$可能不是7的倍数，所以最后一小组的元素可能不足7个），所以共去掉2个小组。

所以在最坏情况下，算法要对不超过$\frac{5}{7}n+8$的元素进行递归，所以算法的最坏情况时间复杂度满足：

%--
\begin{align*}
    W(n) \leq W\bkt{ \ceil{\frac{n}{7}}} + W\bkt{ \frac{5}{7}n+8} + O(n)
\end{align*}

\noindent 同样通过“guess and prove”的方式，我们不难证明$W(n) = O(n)$，具体证明留作习题。

上述证明表明了我们的猜测是准确的，将分组变得更细（从5个一组变成7个一组）并不会使算法的时间复杂度变坏。
但是在大O的意义下，分7个一组也并不会对算法的时间复杂度带来本质改进，而只会让算法的具体实现变得更复杂。
这是因为选择问题显然具有$\Omega(n)$的下界（参见\ref{Sec:Median-Lower-Bound}），而$O(n)$的\alg{Selection-WLinear}算法已经是最优的。
所以实际应用中，当我们需要确保最坏情况时间复杂度时，我们主要使用5个元素一组的\alg{Selection-WLinear}算法。